\section*{الفصل \ref{ch:b3:computational-chemistry} --- الكيمياء الحاسوبية: هارتري--فوك ونظرية دالية الكثافة}

\begin{solution}{exo:b3:computational-chemistry:1}
$\dd E/\dd\alpha = \frac32 - \sqrt{2/\pi\alpha} = 0$ تعطي $\alpha = 8/9\pi = 0.283\,
a_0^{-2}$ و$E = \frac32\alpha - 2\sqrt{2\alpha/\pi} = -4/3\pi = \qty{-0.4244}{\hartree}$،
أي أعلى بنسبة 15\,\% من القيمة الدقيقة $-0.5$: فليس للدالة الغاوسية الواحدة الرأس المدبّب ولا
ذيل الدالة $1s$.
\end{solution}

\begin{solution}{exo:b3:computational-chemistry:2}
STO-3G: خمس دوال لكل C ($1s$ و$2s$ وثلاث $2p$) ودالة واحدة لكل H: $6 \times 5 + 6 =
36$. 6-31G(d): لكل C، 1 (داخلية) + $2 \times 4$ (تكافؤ مجزأ) + 6 ($d$) $= 15$؛
ولكل H، 2: $6 \times 15 + 6 \times 2 = 102$.
\end{solution}

\begin{solution}{exo:b3:computational-chemistry:3}
\omterm{def:b3:quantum-model-systems:schrodinger}{مؤثر هاملتون} الإلكتروني لا يحتوي على كتل النوى: ففي إطار
\omterm{def:b3:computational-chemistry:born-oppenheimer}{تقريب بورن--أوبنهايمر} يكون \omterm{def:b3:computational-chemistry:pes}{سطح الطاقة الكامنة}، ومنه $r_e$ و$k = U''(r_e)$، هو
نفسه. أما $\tilde\omega = \sqrt{k/\mu}/2\pi c$ فيتعلق \omterm{def:b3:quantum-model-systems:oscillator}{بالكتلة المختزلة}، التي
تتضاعف: النسبة $\sqrt2 = 1.414$ (المقيسة 1.413).
\end{solution}

\begin{solution}{exo:b3:computational-chemistry:4}
ثماني ذرات تعطي $3N - 6 = 18$ اهتزازًا: 17 حقيقيًا وواحدًا تخيليًا. \omterm{def:b3:quantum-model-systems:operator}{وقيمة ذاتية}
سالبة واحدة لمصفوفة هيس: نقطة سرج من الرتبة الأولى، أي بنية
انتقالية (هنا مثلًا لدوران أو لحذف).
\end{solution}

\begin{solution}{exo:b3:computational-chemistry:5}
$\zeta = 27/16$، $E = -(27/16)^2 = \qty{-2.8477}{\hartree} = \qty{-77.49}{eV}$.
الخطأ $-2.8477 + 2.9034 = \qty{0.0557}{\hartree} = \qty{1.52}{eV}$، أي
1.9\,\% من الطاقة الكلية --- لكنه أكبر من طاقات تفاعلات كثيرة.
\end{solution}

\begin{solution}{exo:b3:computational-chemistry:6}
$(-1 - E)(-0.5 - E) - (-0.2 - 0.3E)^2 = 0$، أي $0.91E^2 + 1.38E + 0.46 = 0$:
$E = \qty{-1.0217}{\hartree}$ و\qty{-0.4947}{\hartree}. نعم: مزج الدالة
الثانية يخفض الطاقة إلى ما دون $H_{11}$، كما يسمح \omterm{thm:b3:computational-chemistry:variational}{المبدأ
التغايري}.
\end{solution}

\begin{solution}{exo:b3:computational-chemistry:7}
$(102/36)^4 = 64$: نحو 640 دقيقة، أي قرابة 11 ساعة.
\end{solution}

\begin{solution}{exo:b3:computational-chemistry:8}
$0.5782 \times 27.211 = \qty{15.73}{eV}$، أي أعلى بمقدار 0.30 eV من القيمة المقيسة. المدارات
المجمّدة: الأيون الحقيقي يرتخي فتكون طاقته أدنى، فتأتي قيمة كوبمانز مرتفعة جدًا.
الترابط المفقود: للجزيء المتعادل (إلكترونان) طاقة ترابط
أكبر من الأيون (إلكترون واحد)، وهذا يجعل القيمة الحقيقية أكبر.
وهنا يغلب الخطأ الأول.
\end{solution}

\begin{solution}{exo:b3:computational-chemistry:9}
$-0.5960 - 2(-0.4666) = \qty{0.3372}{\hartree} = \qty{9.18}{eV}$ فوق ذرتين.
فالدالة RHF تُبقي 50\,\% من \ce{H+ H-}، وفصل بروتون عن
هيدريد يكلّف الفرق بين طاقة تأين H 
والألفة الإلكترونية للذرة H، وهو فرق تدخله الطاقة في متوسطها.
\end{solution}

\begin{solution}{exo:b3:computational-chemistry:10}
$E(c_1^2 + 2c_1c_2S + c_2^2) = c_1^2H_{11} + 2c_1c_2H_{12} + c_2^2H_{22}$.
بالاشتقاق بالنسبة إلى $c_1$ مع $\partial E/\partial c_1 = 0$:
$E(2c_1 + 2c_2S) = 2c_1H_{11} + 2c_2H_{12}$، أي $(H_{11} - E)c_1 + (H_{12} -
ES)c_2 = 0$؛ وبالمثل $(H_{12} - ES)c_1 + (H_{22} - E)c_2 = 0$.
\end{solution}

\begin{solution}{exo:b3:computational-chemistry:11}
$k \approx [E(1.30) - 2E(1.35) + E(1.40)]/(0.05)^2 = 0.001417/0.0025 =
\qty{0.567}{\hartree\per\bohr\squared} = \qty{882}{N/m}$. ومع $\mu = m_H/2 =
\qty{8.37e-28}{kg}$، يكون $\omega = \sqrt{k/\mu} = \qty{1.03e15}{s^{-1}}$ 
و$\tilde\omega = \omega/2\pi c = \qty{5452}{cm^{-1}}$، أي أعلى بنسبة 24\,\% من القيمة المقيسة
$\tilde\omega_e$: فمنحنى RHF في الأساس الأدنى شديد الانحدار (لا ترابط، وأساس
صغير جدًا)، ولهذا تُخفَّض الترددات المحسوبة بعامل.
\end{solution}

\begin{solution}{exo:b3:computational-chemistry:12}
الفروق التي لا تتجاوز بضعة \unit{kJ/mol} وتتضمن رابطة هيدروجينية تحتاج إلى الترابط 
والتشتت: دالية هجينة مع تصحيح للتشتت، أو الأفضل طريقة
دوال موجية مترابطة للطاقات النهائية، مع أساس مستقطَب
ثلاثي التجزئة يتضمن \omterm{def:b3:computational-chemistry:split-valence}{دوالًا منتشرة}. حسّن هندسة المتشكّلين كليهما بالمستوى
نفسه، وتأكد أنهما قيمتان دنييان بالترددات (كلها حقيقية)، وأضف طاقات
النقطة الصفرية، واختبر الأساس بحساب واحد أكبر، وقارن بنظام
قريب طاقته التشكّلية معروفة.
\end{solution}

\begin{solution}{pb:b3:computational-chemistry:1}
\textbf{1.} يُستبدل بكل \omterm{def:b3:computational-chemistry:basis}{مدار من نمط سلاتر} في أساس أدنى تقليصٌ
ثابت من ثلاث دوال غاوسية مضبوطة عليه.
\textbf{2.} نواة الهيليوم تجذب إلكتروناتها أكثر، فتكون دالتها $1s$
أكثر تراصًا (أسس أكبر). $6.3624/3.4253 = 1.857 = (1.69/1.24)^2$.
\textbf{3.} دالتا أساس، وإلكترونان، ومدار مشغول واحد (ومدار
شاغر واحد).
\textbf{4.} $V_{\mathrm{nn}} = Z_{\mathrm{He}}Z_{\mathrm H}/R = 2/1.4632 =
\qty{1.3669}{\hartree}$.
\textbf{5.} الدالتان متداخلتان بقوة (54\,\%): الذرتان قريبتان
بما يكفي لتترابطا.
\textbf{6.} $h_{11}$ طاقة إلكترون في دالة He بوجود النواتين
دون أي إلكترون آخر؛ ونواة He (الشحنة 2) تمسكه بإحكام
أكبر بكثير.
\textbf{7.} $(h_{11} - \varepsilon)(h_{22} - \varepsilon) - (h_{12} - \varepsilon S)^2 = 0$:
$0.71185\varepsilon^2 + 2.79292\varepsilon + 2.44969 = 0$.
\textbf{8.} $\varepsilon = \qty{-2.600}{\hartree}$ و\qty{-1.324}{\hartree}.
\textbf{9.} تُهمل $\mathbf h$ التنافر بين الإلكترونين؛ فيأتي مدار
القلب شديد التقلص وشديد الانخفاض.
\textbf{10.} $F_{\mu\nu} = h_{\mu\nu} + \sum_{\lambda\sigma}P_{\lambda\sigma}[(\mu\nu|\sigma\lambda) -
\frac12(\mu\lambda|\sigma\nu)]$، مع $P_{\lambda\sigma} = 2C_{\lambda1}C_{\sigma1}$.
\textbf{11.} التنافر الكولومي بين كل إلكترون وسحابة شحنة
الآخر (والتبادل، في الحالة العامة).
\textbf{12.} تتعلق $\mathbf F$ بالمصفوفة $\mathbf P$، التي تتعلق بالمدارات المستخرجة
من $\mathbf F$: فالمعادلات غير خطية وتُحل بتقريبات
متتالية حتى الاتساق الذاتي.
\textbf{13.} $c_1^2 + c_2^2 + 2c_1c_2S = 0.7684 + 0.0410 + 0.1906 = 1.0000$.
\textbf{14.} He: $1.5369 + 0.1906 = 1.727$؛ H: $0.0820 + 0.1906 = 0.273$
إلكترون.
\textbf{15.} الشحنتان: He $+0.27$، وH $+0.73$. فالشحنة الموجبة تقع في معظمها على
الهيدروجين: \ce{HeH+} بروتون مرتبط بذرة هيليوم، \ce{He + H+}، كما هو متوقع
لأن طاقة تأين He (24.6 eV) تفوق كثيرًا طاقة تأين H (13.6 eV).
\textbf{16.} $-\varepsilon_1 = \qty{1.633}{\hartree} = \qty{44.4}{eV}$.
\textbf{17.} $E_{\mathrm{el}} = -2.8418 - 1.3669 = \qty{-4.2087}{\hartree}$.
\textbf{18.} تبلغ ثلاث دورات \qty{-2.8418}{\hartree}. كل دورة تعطي
محددًا، طاقته حد أعلى لطاقة هارتري--فوك المتقاربة
(\omterm{thm:b3:computational-chemistry:variational}{المبدأ التغايري})، والتكرارات تحسّنه.
\textbf{19.} الأساس ذو $\zeta = 2.0925$ يعطي الطاقة الأدنى، فهو
الأفضل لهذا الجزيء (\omterm{thm:b3:computational-chemistry:variational}{المبدأ التغايري}): دالة الهيليوم
أكثر تراصًا في \ce{HeH+} منها في الذرة الحرة، التي ينسخ الأساس المعياري
قيمتها $\zeta = 1.69$.
\textbf{20.} المدار المضاد للربط الفارغ $\sigma^*$؛ وطاقته تقارب
سالب الألفة الإلكترونية للأيون \ce{HeH+} (تقريبًا رديئًا، في أساس صغير كهذا).
\textbf{21.} صفر (لا إلكترون)؛ $\qty{-2.8078}{\hartree}$.
\textbf{22.} $-2.8078 - (-2.8418) = \qty{0.0340}{\hartree} = \qty{89}{kJ/mol}$.
\textbf{23.} نحو نصف القيمة المقيسة: فالأساس الأدنى لا يستطيع أن يستقطب
كثافة الهيليوم نحو البروتون (لا دوال $p$ فيه)، فتأتي الرابطة
ضعيفة جدًا. (وتصحيحات النقطة الصفرية و\qty{298}{K} صغيرة بجانب هذا.)
\textbf{24.} $-(24.5874 + 54.4178)/27.211 = \qty{-2.9034}{\hartree}$؛ وذرة STO-3G
أعلى بمقدار \qty{0.0956}{\hartree} (\qty{2.6}{eV}).
\textbf{25.} تتعلق الهندسات بكيفية تغيّر الطاقة مع $R$؛ ومعظم
الخطأ، المتركز في القلب، يكاد يكون هو نفسه عند كل هندسة
فيُلغي بعضه بعضًا.
\textbf{26.} \textbf{$E(\ce{HeH+}) = \qty{-2.8418}{\hartree}$ بمستوى RHF/STO-3G (الأساس
المعياري)، عند \qty{1.4632}{\bohr}} (و\qty{-2.8607}{\hartree} مع القيمة الكلاسيكية
$\zeta_{\mathrm{He}} = 2.0925$).
\end{solution}
